\section{Obligations and existing foundations}
The table retains the detailed mechanism-to-obligation mapping. Bibliographic
entries for the named mechanisms are provided in the main manuscript's related
work. The mapping identifies reused foundations and the role of each obligation
in this controlled comparison; it does not assert that existing systems cannot
express the same relations.

\begin{table}[tbp]
\caption{Existing foundations and their role in the corrected-execution comparison. P1--P6 are defined in main-text Section 3.2.}\label{tab:obligation-positioning}
\centering\small
\begin{tabularx}{\textwidth}{@{}>{\raggedright\arraybackslash}p{0.12\textwidth}>{\raggedright\arraybackslash}p{0.37\textwidth}>{\raggedright\arraybackslash}X@{}}
\toprule
Obligation & Existing basis & Role in this study \\
\midrule
P1 & Continuity Kernel binds proposal and sealed candidate; ReSA tracks the consumed object. & Select context or reviewed-instance agreement as the binding condition while holding correction and shared checks fixed; observe continuation. \\
P2 & ReSA requires preserved or atomically validated consumed versions. & Check internal record, field and evidence consistency against the submitted immutable candidate. Cross-candidate equality belongs to P1; external referents need a separate consumption check. \\
P3 & Guarded Commits binds authority and pre-edit, delta and post-edit hashes. & Check principal authority and the explicit authorized value, retaining the machine proposal separately from the correction and committed value. \\
P4 & Optimistic concurrency control and Continuity Kernel validate predecessors. & Keep freshness shared across policies, separate from instance mismatch. \\
P5 & Transactions and Continuity Kernel install an atomic accepted unit. & Install the complete corrected successor, sources, transitions and grant consumption together. \\
P6 & Provenance and curated-database copying retain derivation and source relations. & Link changed fields to authorized transitions and copy unchanged fields' exact prior sources; check these relations in persistence. \\
\bottomrule
\end{tabularx}
\end{table}

P1 is policy-relative: context binding requires $K_{\mathrm{eq}}$, and bound
binding adds $I_{\mathrm{eq}}$. P2 is context/evidence integrity and P3 is
principal/authorized-value integrity in both policies. The source-level Alloy
names use a historical numbering scheme; Section~\ref{app:formal-mapping}
gives their precise correspondence and verification limits.

The comparison's admission-set inclusions follow from the declared predicates.
The store checks assess implementation of the shared obligations; the agent
trajectories assess choices and costs after admission or rejection. The
observed reuse, renewed authorization and confirmation actions cannot be
derived from the binding predicates alone.
